\chapter{Grassmann Numbers and Weyl Spinors}\label{app:grassmann}

By the anticommuting nature of Grassmann numbers, the square of any of $\theta^A$ or $\widebar\theta_{\dot{A}}$ are zero.
This has the very useful implication that any function of Grassmann coordinates can be written as
\begin{equation}
    f(\theta_A) = c_0 + c_1\theta_A.
\end{equation}
In addition, we can do calculus on Grassmann numbers, defining the derivative and integral with respect to Grassmann numbers by
\begin{subequations}
\begin{gather}
    \partial_A \theta^B \equiv \frac{\partial \theta^B}{\partial\theta^A} = \delta^B_A,
    \\
    \int d\theta_A \equiv 0, \quad \int d\theta_A \theta_A \equiv 1.
\end{gather}
\end{subequations}
In addition, we normalize the integrals over contracted Grassmann coordinates, such that
\begin{equation}
    \int d^2\theta \; \thetatwo = \int d^2\widebar\theta \; \thetabartwo = \int d^4\theta \; \thetafour = 1,
\end{equation}
where we have defined Weyl spinor contractions such that for two Weyl spinors $\psi$ and $\chi$ we have
\begin{subequations}
\begin{align}
    \psi\chi &\equiv \psi^A\chi_A = \psi_2\chi_1-\psi_1\chi_2,
    \\
    \widebar\psi\widebar\chi &\equiv \widebar\psi_{\dot{A}}\widebar\chi^{\dot{A}} = \widebar\psi_{\dot{1}}\widebar\chi_{\dot{2}} - \widebar\psi_{\dot{2}}\widebar\chi_{\dot{1}}.
\end{align}
\end{subequations}